\documentclass{article}
\usepackage[utf8]{inputenc}
\usepackage[T1]{fontenc}
\usepackage{geometry}
\usepackage{graphicx}
\usepackage{float}
\usepackage{hyperref}
\usepackage[english,indonesian]{babel}

\title{Optimasi Pemrosesan Template LaTeX Menggunakan Pendekatan Berbasis AI}
\author{Academic Support Agent}
\date{\today}

\begin{document}
\maketitle

\begin{abstract}

\end{abstract}

\chapter{PENDAHULUAN}


\chapter{TINJAUAN PUSTAKA}


% Original content (commented out):
% %%%%%%預先準備（不需修改）%%%%%%
\documentclass[12pt]{article}
\input{preamble}

%%%%%%%%自訂設置（請填入）%%%%%%%%
\centerhead{LaTeX 生成範例} %中上方頁首，請填入小論文主題名稱

%%%%%%%%文件內容開始%%%%%%%%%%%%
\begin{document}
\coverpage
    {編譯測試類} %eg.工程技術類
    {LaTeX 生成範例} %填入主題名稱
    %填入作者。若作者只有 1 位 或 2 位，則將 {} 內留白即可
    {01 王小明。國立虛擬高級中學。高三 6 班} %第一作者
    {} %第二作者
    {} %第三作者
    {時間就是最好的老師} %指導老師
    
\h{前言}

    README 可參閱 \url{https://junqi.netlify.app/latexsmallessaytemplate}\\
    或 \url{https://hackmd.io/@coke5151/rJz_LM3cc}
    
    影音教學可參考 \url{https://youtu.be/3Q4QSGgzqMs}

\h{文獻探討}

\h{研究方法}

\h{研究分析與結果}

\h{研究結論與建議}

\clearpage %後方的內容獨立從新的一頁開始
\h{參考文獻}

%%%%%%%%%%%%%%%%%%%%%%%%%%%%%%%%%%%%%%%%%%%%%
%此區塊為參考文獻，請直接編輯 ref.tex 檔案%
\begin{enumerate}[(1)]
\input{ref}
\end{enumerate}
%%%%%%%%%%%%%%%%%%%%%%%%%%%%%%%%%%%%%%%%%%%%%

\end{document}


\end{document}